\section{Weekly Summary}

\begin{tcolorbox}[colback=yellow!10,colframe=black!50,title=AI-Generated Content]
\noindent The summary below was written by Claude (Anthropic) from that week's regression outputs and series descriptions. It has not been reviewed or fact-checked by a human, and should be read as an automated, informal recap -- not analysis.
\end{tcolorbox}

Welcome back to the show where we yank two completely unrelated economic time series out of a hat each day, force them to arm-wrestle in a scatterplot, and then pretend to be shocked by the outcome. This week's roster featured such natural dance partners as Great Lakes wood product GDP paired with European cell phone subscriptions, import prices for beverages and tobacco squaring off against unfilled orders for photographic equipment, and -- my personal favorite -- a Producer Price Index for computer training services introduced to the flow of securities at domestic finance companies. Truly the stuff of economics textbooks. Usual disclaimer, delivered with the enthusiasm of a flight attendant on their ninth leg of the day: none of these pairings mean anything, there is no causal story, and spurious correlation is not a bug here, it is the entire product.

Let's start with Monday's headline act, the wood products versus cell phones showdown. The raw regression gave a polite shrug (p-value around 0.14, nothing to write home about), but -- plot twist nobody ordered -- the DETRENDED version actually came out significant at the 0.02 level. So once you strip the shared march of time out of both series, Michigan sawmills and Eurasian phone plans still seem to wiggle together a little. Is this a genuine relationship? Almost certainly not, unless you think lumberjacks are furiously upgrading to 5G. But credit where it's due: it survived the detrending guillotine, which is more than I can say for our next contestant.

Ah yes, Tuesday. Behold the mighty R-squared of 0.886 between labor-income share and the number of Black Americans not in the labor force -- a p-value with nineteen zeros after the decimal point, the kind of number that makes amateur analysts start drafting Nobel acceptance speeches. And then we detrend, and the R-squared face-plants from 0.89 all the way down to 0.23. This, friends, is the textbook trend mirage: two series drifting across the decades, holding hands only because time was passing for both of them. It still clings to a weak significance after detrending, but that spectacular raw number was mostly just the gravitational pull of history. Wednesday's intellectual-property imports versus durable-goods share pulled the exact same stunt -- gorgeous raw p-value, then a detrended p-value of 0.44 that waved goodbye and left the building entirely.

Rounding out the week: Thursday's government imputations versus official reserve assets managed to be impressively nothing in both regressions (R-squared of 0.0009, a rounding error with ambitions), and Friday served up Mountain Census unemployment against Southwest gasoline spending, which actually held onto significance after detrending (p around 0.037) -- so if you squint, maybe people drive less when the mountains are out of work. Or maybe it's noise wearing a trench coat. As always, these are random pairings tortured into charts for fun, no rigorous claims are being made, and the only thing you should confidently conclude is that I need a hobby that produces fewer footnotes. See you next week.
